% This file is generated from the corresponding Obsidian note.
% Source: Teoricos/GS - Teo24.md
\chapter{Orientación}
\label{ch:teo-24}
\section{Orientacion de espacios vectoriales}
\begin{bookdefinition}{Orientacion de un espacio vectorial}
Sea $V$ un espacio vectorial real de dimension finita. Se define una relacion de equivalencia en el conjunto de bases ordenadas de $V$ como:

Si $\mathcal B_1$ y $\mathcal B_2$ son bases ordenadas, decimos que

\[
\mathcal B_1\sim\mathcal B_2
\]
si el determinante de la matriz de cambio de base de $\mathcal B_1$ a $\mathcal B_2$ es positivo.

Las clases de equivalencia se llaman orientaciones de $V$.

\end{bookdefinition}
\begin{bookexercise}{}
Verificar que la anterior relacion es una relacion de equivalencia en el conjunto de bases ordenadas.

\end{bookexercise}
\begin{bookproof}{}
1.

\bookqed
\end{bookproof}
\begin{bookproposition}{}
La relacion de equivalencia en el conjunto de bases ordenadas de $V$ tiene solo dos clases de equivalencia.

\end{bookproposition}
\begin{bookproof}{}
\begin{enumerate}
\item \textbf{Hay al menos dos clases:} Sea $\mathcal B=\{v_1,\ldots,v_n\}$ una base de $V$. Entonces
\end{enumerate}
\[
\widetilde{\mathcal B}=\{-v_1,v_2,\ldots,v_n\}
\]
tambien es base de $V$, y la matriz de cambio de base tiene determinante

\[
\det\begin{pmatrix}-1& &0\\&\ddots&\\0& &1\end{pmatrix}=-1<0.
\]
\begin{enumerate}[start=2]
\item Por lo tanto $\mathcal B$ y $\widetilde{\mathcal B}$ no estan en la misma clase.
\item \textbf{Hay exactamente dos:} Sean $\mathcal B$ y $\widehat{\mathcal B}$ dos bases ordenadas que no estan en la misma clase. Sea $\mathcal B'$ una tercera base ordenada.
\item Si $\mathcal B'\sim\mathcal B$, listo. Si no, entonces el determinante de la matriz de cambio de base de $\mathcal B$ a $\mathcal B'$ es negativo.
\item Como $\mathcal B$ y $\widehat{\mathcal B}$ no estan en la misma clase, tambien
\end{enumerate}
\[
\det\left(M_{\mathcal B}^{\widehat{\mathcal B}}\right)<0.
\]
\begin{enumerate}[start=6]
\item Sabemos que
\end{enumerate}
\[
M_{\mathcal B}^{\mathcal B'}=M_{\widehat{\mathcal B}}^{\mathcal B'}M_{\mathcal B}^{\widehat{\mathcal B}},
\]
y por tanto

\[
\det\left(M_{\mathcal B}^{\mathcal B'}\right) =\det\left(M_{\widehat{\mathcal B}}^{\mathcal B'}\right)\det\left(M_{\mathcal B}^{\widehat{\mathcal B}}\right).
\]
\begin{enumerate}[start=7]
\item Como los dos determinantes de los extremos tienen signo negativo, se sigue que
\end{enumerate}
\[
\det\left(M_{\widehat{\mathcal B}}^{\mathcal B'}\right)>0,
\]
y entonces $\mathcal B'\sim\widehat{\mathcal B}$.

\bookqed
\end{bookproof}
\begin{bookdefinition}{Espacio vectorial orientado}
Un espacio vectorial orientado es un par $(V,[\mathcal B])$, donde $V$ es un espacio vectorial real de dimension finita y $[\mathcal B]$ es una clase de equivalencia de la relacion de orientacion en el conjunto de bases ordenadas de $V$.

Una base ordenada $\mathcal B$ de $V$ se dice positivamente orientada si $\mathcal B\in[\mathcal B]$.

\end{bookdefinition}
\begin{bookremark}{}
Si $V$ tiene dimension $n$, podemos usar una $n$-forma no nula para dar una orientacion a $V$.

\end{bookremark}
\begin{booklemma}{}
Sea $\omega\in\Lambda^n(V^*)$ y sea $\mathcal B=\{e_1,\ldots,e_n\}$ una base de $V$. Entonces, dados $v_1,\ldots,v_n\in V$ arbitrarios, se cumple que

\[
\omega(v_1,\ldots,v_n)=\det A\ \omega(e_1,\ldots,e_n),
\]
donde $A$ es la matriz cuya columna $j$ son las coordenadas de $v_j$ en la base $\mathcal B$.

\end{booklemma}
\begin{bookproof}{}
\begin{enumerate}
\item Sea $\mathcal B^*=\{e_1^*,\ldots,e_n^*\}$ la base dual de $V^*$. Sabemos que
\end{enumerate}
\[
e_1^*\wedge\cdots\wedge e_n^*
\]
es base de $\Lambda^n(V^*)$. 2. Por lo tanto existe $\lambda\in\mathbb R$ tal que

\[
\omega=\lambda e_1^*\wedge\cdots\wedge e_n^*.
\]
\begin{enumerate}[start=3]
\item Evaluando en la base,
\end{enumerate}
\[
\omega(e_1,\ldots,e_n)=\lambda.
\]
\begin{enumerate}[start=4]
\item Por otra parte, si $A=(a_{ij})$ es la matriz de coordenadas de los $v_j$, entonces
\end{enumerate}
\[
\omega(v_1,\ldots,v_n)=\lambda\det(a_{ij})=\det A\ \omega(e_1,\ldots,e_n).
\]
recordando que $e_{j}^{*}(v_{i})$ nos da la $j$-esima coordenadas de el vector $v_{i}$ escrito en la base $\mathcal{B}$

\bookqed
\end{bookproof}
\phantomsection\label{blk:543d91}
\begin{bookproposition}{}
Sea $V$ un espacio vectorial real de dimension $n$. Cada elemento no nulo $\omega\in\Lambda^n(V^*)$ determina una orientacion $\mathcal O$ de $V$ de la siguiente manera:

\[
\mathcal O=\{(E_1,\ldots,E_n)\text{ base ordenadas}:\omega(E_1,\ldots,E_n)>0\}.
\]
Dos $n$-tensores alternantes no nulos determinan la misma orientacion si y solo si uno es multiplo positivo del otro.

\end{bookproposition}
\begin{bookproof}{}
\begin{enumerate}
\item Sea $\mathcal B=\{e_1,\ldots,e_n\}$ una base ordenada y sea $\omega\in\Lambda^n(V^*)$ no nula. Vamos a ver que la clase de equivalencia de esta base está dada por las bases ordenadas $\{v_1,\ldots,v_n\}$ tales que $\omega(v_1,\ldots,v_n)$ tiene el mismo signo que $\omega(e_1,\ldots,e_n)$
\item Supongamos tenemos otra base ordenada $\{ v_{1},\ldots,v_{n} \}$ entonces:
\end{enumerate}
\[
\begin{aligned}\{ v_{1},\ldots v_{n} \}\in [\{ e_{1},\ldots e_{n} \}]& \iff \det(M_{\widetilde{B}}^{B} )>0\\& \iff \frac{\omega(v_{1},\ldots,v_{n})}{\omega(e_{1},\ldots,e_{n})}=\det(M_{\widetilde{\mathcal B}}^{B} )>0\\&\iff \omega(v_{1},\ldots,v_{n})\ \land \ \omega(e_{1},\ldots,e_{n})\ \text{tienen el mismo signo}\end{aligned}
\]
el primer paso por definicion el segundo por \hyperref[blk:543d91]{Lema} 3. Por lo tanto una orientacion de $V$ esta dada por las bases ordenadas $\{ v_{1},\ldots v_{n} \}$ tales que $\omega(v_{1},\ldots,v_{n})$ tienen el mismo signo. 4. Normalmente llamamos orientacion (denotada $\mathcal{O}_{\omega}$) a las bases tales que $\omega(v_{1},\ldots,v_{n})>0$ osea tienen el mismo signo y son positivas. 5. Si $\widetilde{\omega}$ es otra $n$-forma no nula en $\Lambda^{*}(V^{*})$, claramente

\[
\mathcal{O}_{\omega}=\mathcal{O}_{\widetilde{\omega}}\iff \widetilde{\omega}=\lambda\omega \  \text{ con } \ \lambda>0
\]
\bookqed
\end{bookproof}
\section{Orientacion de variedades}
\begin{bookdefinition}{Variedad orientable y atlas orientado}
Una variedad suave conexa $M$ de dimension $m$ se dice orientable si existe un atlas suave $\mathcal A$ tal que para cualquier par de cartas suaves

\[
(U,\varphi=(x_1,\ldots,x_m)),\qquad (V,\psi=(y_1,\ldots,y_m))
\]
en $\mathcal A$, con $U\cap V\ne\varnothing$, se cumple

\[
\det\left(\frac{\partial y_j}{\partial x_i}\right)>0\qquad \text{sobre }U\cap V.
\]
Un atlas con esta propiedad se llama \textbf{atlas orientado}.

\end{bookdefinition}
\begin{bookremark}{}
Dicho \textbf{atlas orientado} induce una orientacion en cada $T_{p}M$ con $p\in M$ dada por la clase de equivalencia de

\[
\left\{ \frac{\partial}{\partial x_{1}}\bigg|_{p},\ldots, \frac{\partial}{\partial x_{m}}\bigg|_{p}\right \}
\]
para cualquier carta $(U,\varphi=(x_{1},\ldots,x_{m}))\in \mathcal{A}$ Esto es por que notamos que como el \textbf{atlas orientado} tenemos que

\[
0<\det\left(\frac{\partial y_j}{\partial x_i}\right)_{ij}=\det(J(\psi\circ\varphi ^{-1}))\circ\varphi \qquad\text{sobre}\quad U\cap V
\]
En particular si la otra carta que tomamos $(V,\psi=y_{1},\ldots y_{m})\in \mathcal A$ cumple $p\in V$ entonces por definicion de \textbf{atlas orientado} tenemos que

\[
\det \big[ J(\psi\circ\varphi ^{-1})(\varphi(p))\big]>0
\]
pero esta es la matriz de cambio de base entre

\[
\left\{\frac{\partial}{\partial x_1}\bigg|_p,\ldots,\frac{\partial}{\partial x_m}\bigg|_p\right\}\qquad\text{y}\qquad \left\{\frac{\partial}{\partial y_1}\bigg|_p,\ldots,\frac{\partial}{\partial y_m}\bigg|_p\right\}.
\]
con lo cual en el espacio vectorial de $T_{p}M$ estas bases pertenecen a la misma clase de equivalencia.

\end{bookremark}
\begin{bookexample}{}
\begin{enumerate}
\item $\mathbb R^n$ es orientable.
\item $S^n$ es orientable con $n\geq 2$
\end{enumerate}
\end{bookexample}
\begin{bookproposition}{}
En general, cualquier variedad que pueda cubrirse con dos cartas suaves tal que $U\cap V$ sea conexo

\end{bookproposition}
\begin{bookproof}{}
\begin{enumerate}
\item Sea $M$ una variedad que puede cubrirse con dos cartas suaves $(U,\varphi=(x^1,\ldots,x^m))$ y $(V,\psi=(y^1,\ldots,y^m))$ tales que $U\cap V$ sea conexo.
\item Tenemos la función
\end{enumerate}
\[
f:U\cap V\to\mathbb R \qquad q\mapsto\det\!\left(\frac{\partial y^j}{\partial x^i}(q)\right)
\]
\begin{enumerate}[start=3]
\item Como $\left(\frac{\partial y^j}{\partial x^i}(q)\right)$ es una matriz de cambio de base, se tiene $f(q)\neq0$ para todo $q\in U\cap V$.
\item Además, $f$ es continua. Como $U\cap V$ es conexo, $f(U\cap V)$ es un conjunto conexo de $\mathbb R$ que no contiene al $0$. Por lo tanto, o bien $f>0$ en todo $U\cap V$, o bien $f<0$ en todo $U\cap V$.
\item En el primer caso, $M$ es orientable por definición.
\item En el segundo caso, cambiamos una de las cartas y definimos $\tilde\psi=(-y^1,y^2,\ldots,y^m)$
\item Entonces
\end{enumerate}
\[
\det\!\left(\frac{\partial\tilde y^j}{\partial x^i}\right)=-\det\!\left(\frac{\partial y^j}{\partial x^i}\right)>0
\]
para todo punto de $U\cap V$. 8. Por consiguiente, el atlas $\mathcal A=\{(U,\varphi),(V,\tilde\psi)\}$ es una orientación de $M$. 9. Luego, $M$ es orientable.

\bookqed
\end{bookproof}
\begin{bookdefinition}{Carta positivamente/negativamente orientada}
Sea $(M, \mathcal{A})$ una variedad orientada y sea $(U,\varphi=(x_1,\ldots,x_m))$ una carta suave arbitraria de $M$. Se dice que $(U,\varphi)$ esta positivamente orientada si la base ordenada

\[
\left\{\frac{\partial}{\partial x_1}\bigg|_p,\ldots,\frac{\partial}{\partial x_m}\bigg|_p\right\}
\]
tiene orientacion inducida por $\mathcal A$ para todo $p\in U$ Es decir para cada $p\in U$ sucede que

\[
\det\left(\frac{\partial y_j}{\partial x_i}\bigg|_{p}\right)>0
\]
para cualquier carta $(V,\psi=(y_{1},\ldots y_{n})$ en $\mathcal A$ tal que $p\in U\cap V\neq 0$ Analogamente se define carta negativamente orientada.

\end{bookdefinition}
\begin{bookproposition}{}
Sea $(M,\mathcal A)$ una variedad orientada y sea $(U,\varphi=(x_1,\ldots,x_m))$ una carta suave de $M$ con $U$ conexo. Entonces $(U,\varphi)$ es positivamente orientada o negativamente orientada.

\end{bookproposition}
\begin{bookproof}{}
\begin{enumerate}
\item Consideremos los subconjuntos de $U$
\end{enumerate}
\[
U_+=\left\{p\in U:\left\{\frac{\partial}{\partial x_1}\bigg|_p,\ldots,\frac{\partial}{\partial x_m}\bigg|_p\right\}\text{ es positivamente orientada con respecto a } \mathcal A\right\},
\]
\[
U_-=\left\{p\in U:\left\{\frac{\partial}{\partial x_1}\bigg|_p,\ldots,\frac{\partial}{\partial x_m}\bigg|_p\right\}\text{ es negativamente orientada con respecto a }\mathcal A\right\}.
\]
notar que en estos conjuntos hablamos de orientacion como espacios vectoriales 2. Vale la pena recalcar que $U_{+}$ y $U_{-}$ están bien definidas pues $\mathcal A$ define una orientación en cada $T_pM$, con $p\in M$, que no depende de la carta que se elija en $\mathcal A$. 3. Veamos que ambos conjuntos son abiertos de $U$ 4. \textbf{$U_{+}$ es abierto de $M$}: Sea $p\in U_+$ y sea $(V,\psi=(y_1,\ldots,y_n))\in\mathcal A$ una carta de $p$. 5. Luego $p\in U\cap V$ y por definicion de $U_{+}$ tenemos que

\[
\det\!\left(\frac{\partial y_i}{\partial x_j}(p)\right)>0
\]
\begin{enumerate}[start=6]
\item Pensemos en la función
\end{enumerate}
\[
f:U\cap V\longrightarrow\mathbb R,\qquad q\longmapsto \det\!\left(\frac{\partial y_i}{\partial x_j}(q)\right)
\]
\begin{enumerate}[start=7]
\item Como las funciones
\end{enumerate}
\[
\frac{\partial y_i}{\partial x_j}:U\cap V\to\mathbb R
\]
son $C^\infty$, en particular son continuas, tenemos que $f$ es continua. 8. Y dado que $f(p)>0$, entonces $f>0$ en un abierto $W$ de $p$ de $U\cap V$, el cual también es abierto de $M$. 9. Esto quiere decir que para todo $q\in W$, la base

\[
\left\{\frac{\partial}{\partial x_1}\bigg|_q,\ldots,\frac{\partial}{\partial x_n}\bigg|_q\right\}
\]
está positivamente orientada. 10. Por lo tanto $p\in W\subseteq U_{+}$ con $W$ abierto. Luego $U_{+}$ es abierto. 11. Análogamente, $U_{-}$ es un abierto de $M$. 12. Como $U_{+}$ y $U_{-}$ son disjuntos, y $M$ es conexo, entonces $M=U_{+}$ o $M=U_{-}$, que es lo que queríamos probar.

\bookqed
\end{bookproof}
\begin{bookremark}{}
Si una carta $(U,\varphi)$ no esta positivamente orientada, podemos cambiar una coordenada, por ejemplo

\[
(U,\widetilde\varphi)=(-x_1,x_2,\ldots,x_m),
\]
y obtener una carta positivamente orientada.

\end{bookremark}
\begin{bookproposition}{Caracterizacion de orientabilidad}
Sean $M$ una variedad suave conexa de dimension $n$. Entonces las siguientes afirmaciones son equivalentes:

\begin{itemize}
\item (a) $M$ es orientable.
\item (b) Existe una $n$-forma continua nunca nula sobre $M$.
\item (c) Existe una $n$-forma suave nunca nula sobre $M$.
\end{itemize}
\end{bookproposition}
\begin{bookproof}{}
\begin{itemize}
\item \textbf{$(c)\Rightarrow(b)$} es inmediato.
\item \textbf{$(b)\Rightarrow(a)$}.
\end{itemize}
\begin{enumerate}
\item Tenemos $\omega$ una $n$-forma no nula, recordamos que un $n$-tensor alternante no nulo define una orientación (en espacios vectoriales)
\item Usamos la continuidad de la $n$-forma $\omega$ para definir un atlas de orientación.
\item Arrancamos con $\mathcal A=\{(U_\alpha,\varphi_\alpha=(x_1^\alpha,\ldots,x_n^\alpha))\}_{\alpha\in I}$ un atlas suave de $M$, donde podemos suponer que cada $U_\alpha$ es un conjunto conexo.
\item Recordar que las componentes conexas en un espacio topológico localmente conexo por caminos son abiertas.
\item Para cada $(U_\alpha,\varphi_\alpha)\in\mathcal A$, consideramos la función
\end{enumerate}
\[
f_\alpha:U_\alpha\to\mathbb R\qquad p\mapsto \omega_p\!\left(\left.\frac{\partial}{\partial x_1^\alpha}\right|_p,\ldots,\left.\frac{\partial}{\partial x_n^\alpha}\right|_p\right)
\]
\begin{enumerate}[start=6]
\item Como $\omega$ es nunca nula y continua, entonces $f_\alpha$ es continua y nunca nula.
\item Como $U_\alpha$ es conexo, se sigue que $f_\alpha>0$ o bien $f_\alpha<0$ en todo $U_\alpha$.
\item Si $f_\alpha<0$, construimos la nueva carta suave $(U_\alpha,\widetilde\varphi_\alpha=(-x_1^\alpha,x_2^\alpha,\ldots,x_n^\alpha))$.
\item Nuestro atlas de orientación es
\end{enumerate}
\[
\mathcal A_0=\{(U_\alpha,\varphi_\alpha)\in\mathcal A:f_\alpha>0\}\cup\{(U_\alpha,\widetilde\varphi_\alpha):(U_\alpha,\varphi_\alpha)\in\mathcal A,\ f_\alpha<0\}.
\]
\begin{itemize}
\item \textbf{$\mathcal A_0$ es un atlas suave de $M$}: Sus cartas son suavemente compatibles. Veamoslo:
\end{itemize}
\begin{enumerate}
\item Si $(U_\alpha,\varphi_\alpha),(U_\beta,\varphi_\beta)$ son t.q. $f_\alpha>0$, $f_\beta>0$ y $U_\alpha\cap U_\beta\neq\varnothing$, entonces son suavemente compatibles porque están en $\mathcal A$. Y $\mathcal A$ es atlas suave
\item Si $(U_\alpha,\varphi_\alpha)$ y $(U_\beta,\widetilde\varphi_\beta)$ donde $f_\alpha>0$, $f_\beta<0$ y $U_\alpha\cap U_\beta\neq\varnothing$. Consideremos
\end{enumerate}
\[
T:\mathbb R^n\to\mathbb R^n,\qquad (x^1,\ldots,x^n)\mapsto(-x^1,x^2,\ldots,x^n),
\]
la cual es una transf. lineal t.q. $T^2=\operatorname{Id}_{\mathbb R^n}$ y

\[
\widetilde\varphi_\beta=T\circ\varphi_\beta,
\]
así

\[
\widetilde\varphi_\beta\circ\varphi_\alpha^{-1}=T\circ \underbrace{\varphi_\beta\circ\varphi_\alpha^{-1}}_{\text{suave, pues }\mathcal A\text{ es atlas}}
\]
y

\[
\varphi_\alpha\circ\widetilde\varphi_\beta^{-1}=\varphi_\alpha\circ(T\circ\varphi_\beta)^{-1}=\underbrace{\varphi_\alpha\circ\varphi_\beta^{-1}}_{\text{suave, pues }\mathcal A\text{ es atlas}}\circ T^{-1},
\]
\begin{enumerate}[start=3]
\item Y finalmente si $(U_\alpha,\widetilde\varphi_\alpha)$, $(U_\beta,\widetilde\varphi_\beta)$ son t.q $f_{\alpha }<0$ y $f_{\beta}<0$ entonces son suavemente compatibles:
\end{enumerate}
\[
\widetilde\varphi_\beta\circ\widetilde\varphi_\alpha^{-1}=(T\circ\varphi_\beta)\circ(\varphi_\alpha^{-1}\circ T^{-1})=T\circ(\varphi_\beta\circ\varphi_\alpha^{-1})\circ T^{-1}
\]
\begin{itemize}
\item \textbf{Para todo $p\in M$, las cartas de $\mathcal A_0$ definen la misma orientación en $T_pM$:}
\end{itemize}
\begin{enumerate}
\item Sean $(U,\varphi)$, $(V,\psi)\in\mathcal A_0$ t.q. $U\cap V\neq\varnothing$. Tenemos tres casos
\end{enumerate}
\begin{itemize}
\item \textbf{1er Caso:}

\begin{enumerate}
\item Primero tenemos
\end{enumerate}
\end{itemize}
\[
(U,\varphi)=(U_\alpha,\varphi_\alpha),\qquad (V,\psi)=(U_\beta,\varphi_\beta),\qquad f_\alpha>0,\ f_\beta>0.
\]
entonces

\[
0<f_\alpha=\omega\!\left(\frac{\partial}{\partial x_\alpha^1},\ldots,\frac{\partial}{\partial x_\alpha^n}\right)=\det\!\left(\frac{\partial x_\beta^j}{\partial x_\alpha^i}\right)\omega\!\left(\frac{\partial}{\partial x_\beta^1},\ldots,\frac{\partial}{\partial x_\beta^n}\right)=\det\!\left(\frac{\partial x_\beta^j}{\partial x_\alpha^i}\right)f_\beta.
\]
\begin{Verbatim}[fontsize=\small]
 2. Luego 
\end{Verbatim}
\[
\det\!\left(\frac{\partial x_\beta^j}{\partial x_\alpha^i}\right)>0.
\]
\begin{itemize}
\item \textbf{2do caso:}

\begin{enumerate}
\item Si $(U,\varphi)=(U_\alpha,\varphi_\alpha)$ y $(V,\psi)=(U_\beta,\widetilde\varphi_\beta)$, con $f_\alpha>0$ y $f_\beta<0$
\item Entonces, repitiendo la misma cuenta de arriba, tenemos
\end{enumerate}
\end{itemize}
\[
\begin{aligned} 0<f_\alpha & =\omega\!\left(\frac{\partial}{\partial x_\alpha^1},\ldots,\frac{\partial}{\partial x_\alpha^n}\right)\\ &=\det\!\left(\frac{\partial \widetilde x_\beta^j}{\partial x_\alpha^i}\right)\omega\!\left(\frac{\partial}{\partial \widetilde x_\beta^1},\ldots,\frac{\partial}{\partial \widetilde x_\beta^n}\right)\\&=\det\!\left(\frac{\partial \widetilde x_\beta^j}{\partial x_\alpha^i}\right)\omega\!\left(-\frac{\partial}{\partial x_\beta^1},\ldots,\frac{\partial}{\partial x_\beta^n}\right)\\&=-\det\!\left(\frac{\partial \widetilde x_\beta^j}{\partial x_\alpha^i}\right)\omega\!\left(\frac{\partial}{\partial x_\beta^1},\ldots,\frac{\partial}{\partial x_\beta^n}\right)\\&=-\det\!\left(\frac{\partial \widetilde x_\beta^j}{\partial x_\alpha^i}\right)f_{\beta }\end{aligned}
\]
\begin{Verbatim}[fontsize=\small]
 3. Pero como $0>f_{\beta}$ tenemos que 
\end{Verbatim}
\[
\det\!\left(\frac{\partial \widetilde x_\beta^j}{\partial x_\alpha^i}\right)>0
\]
\begin{itemize}
\item \textbf{3er Caso}

\begin{enumerate}
\item Por último, si $(U,\varphi)=(U_\alpha,\widetilde\varphi_\alpha)$ y $(V,\psi)=(U_\beta,\widetilde\varphi_\beta)$ con $f_\alpha<0$ y $f_\beta<0$,
\item Entonces tenemos
\end{enumerate}
\end{itemize}
\[
\begin{aligned} 0>f_\alpha & =\omega\!\left(\frac{\partial}{\partial x_\alpha^1},\ldots,\frac{\partial}{\partial x_\alpha^n}\right)\\&=-\omega\!\left(\frac{\partial}{\partial \widetilde x_\alpha^1},\ldots,\frac{\partial}{\partial \widetilde x_\alpha^n}\right)\\&=\det\!\left(\frac{\partial \widetilde x_\beta^j}{\partial \widetilde x_\alpha^i}\right)-\omega\!\left(\frac{\partial}{\partial \widetilde x_\beta^1},\ldots,\frac{\partial}{\partial \widetilde x_\beta^n}\right)\\& =\det\!\left(\frac{\partial \widetilde x_\beta^j}{\partial \widetilde x_\alpha^i}\right)\omega\!\left(\frac{\partial}{\partial  x_\beta^1},\ldots,\frac{\partial}{\partial  x_\beta^n}\right)\\&=\det\!\left(\frac{\partial \widetilde x_\beta^j}{\partial \widetilde x_\alpha^i}\right)f_{\beta }\end{aligned}
\]
\begin{Verbatim}[fontsize=\small]
 3. Con lo cual 
\end{Verbatim}
\[
\det\!\left(\frac{\partial \widetilde x_\beta^j}{\partial \widetilde x_\alpha^i}\right)>0
\]
\begin{enumerate}[start=4]
\item Luego todas las funciones de transición de $\mathcal A_0$ tienen determinante positivo. Por consiguiente, todas las cartas de $\mathcal A_0$ determinan la misma orientación.
\end{enumerate}
\begin{itemize}
\item \textbf{$(a)\Rightarrow (c)$}
\end{itemize}
\begin{enumerate}
\item Tenemos el atlas de orientación $\mathcal A$. Queremos construir $\omega\in\Omega^n(M)$.
\item Para cada $(U_\alpha,\varphi_\alpha)\in\mathcal A$, con $\varphi_\alpha=(x_\alpha^1,\ldots,x_\alpha^n)$, consideremos
\end{enumerate}
\[
\xi_\alpha=(dx_\alpha^1)\wedge\cdots\wedge(dx_\alpha^n).
\]
\begin{enumerate}[start=3]
\item Tomemos una partición de la unidad subordinada al cubrimiento $\{U_\alpha\}_{\alpha\in I}$, digamos $\{\rho_\alpha\}_{\alpha\in I}$.
\item Definimos sobre $M$
\end{enumerate}
\[
\rho_\alpha\,\xi_\alpha=\begin{cases}\rho_\alpha(q)\,\xi_\alpha(q),&q\in U_\alpha,\\0,&q\notin U_\alpha.\end{cases}
\]
\begin{enumerate}[start=5]
\item Como $\rho_\alpha$ es suave, usando el criterio de suavidad para formas, tenemos que $\rho_\alpha\,\xi_\alpha$ es suave.
\item Definimos entonces
\end{enumerate}
\[
\omega=\sum_{\alpha\in I}\rho_\alpha\,\xi_\alpha.
\]
\begin{enumerate}[start=7]
\item Y $\omega$ resulta suave, pues $\{\operatorname{supp}\rho_\alpha\}_{\alpha\in I}$ es un conjunto localmente finito luego para cada $p\in M$ existe un entorno abierto $V$ de $p$ tal que sólo un número finito de las formas $\rho_\alpha\xi_\alpha$ son no nulas en $V$. En consecuencia, $\omega|_V$ es una suma finita de formas suaves, y por tanto es suave.
\item Además, para cada $p\in M$, $\omega_{p}$ es $n$-lineal y alternante.
\item Sólo resta verificar que $\omega$ nunca se anula. Para esto recordemos que $\sum_{\alpha\in I}\rho_\alpha(p)=1$ y $\rho_\alpha(p)\ge0$ para todo $\alpha\in I$.
\item Dado $p\in M$ como $\sum_{\alpha\in I}\rho_\alpha(p)=1$ existe $\beta\in I$ tal que $\rho_\beta(p)>0$ osea $p\in \operatorname{supp} \rho_{\alpha }\subseteq U_{\beta}$
\item La carta $(U_\beta,\varphi_\beta=(x_\beta^1,\ldots,x_\beta^n))$ determina la base $\{\frac{\partial}{\partial x_\beta^1}\big|_p,\ldots,\frac{\partial}{\partial x_\beta^n}\big|_p\}$ de $T_pM$.
\item Entonces
\end{enumerate}
\[
\begin{aligned}\omega_{p}\!\left(\frac{\partial}{\partial x_\beta^1}\bigg|_p,\ldots,\frac{\partial}{\partial x_\beta^n}\bigg|_p\right)&= \rho_{\beta }(p)\xi_{\beta}\!\left(\frac{\partial}{\partial x_\beta^1}\bigg|_p,\ldots,\frac{\partial}{\partial x_\beta^n}\bigg|_p\right)+\sum_{\substack{\alpha\in I \\ \alpha  \neq \beta }}^{n} \rho_\alpha(p)\,\xi_{\alpha}\!\left(\frac{\partial}{\partial x_\beta^1}\bigg|_p,\ldots,\frac{\partial}{\partial x_\beta^n}\bigg|_p\right)\\&=\rho_{\beta }(p)+\sum_{\substack{\alpha\in I \\ \alpha  \neq \beta }}^{n} \rho_\alpha(p)\det\left(\frac{\partial x_{\alpha }^{j } }{\partial x_{\beta }^{i}}\right) \end{aligned}
\]
para el segundo igual usamos la definicion de evaluar una $n$-forma y la propiedad $dx(\frac{\partial }{\partial y})=\frac{\partial }{\partial y}x$ 13. Y el resultado de esta suma es no negativo dado que $\rho_{\beta}(p)>0$ y $\rho_{\alpha }\geq 0$ y el $\det$ es mayor que $0$ por que ambas cartas pertencen a $\mathcal A$ el atlas de orientacion

\bookqed
\end{bookproof}
\section{Esto no lo dio en clase}
\begin{booktheorem}{Orientabilidad de hipersuperficies}
Sea $M$ una subvariedad de dimension $n$ en $\mathbb R^{n+1}$. Probar que $M$ es orientable si y solo si $M$ admite un campo normal diferenciable nunca nulo; es decir, una funcion diferenciable

\[
\xi:M\to\mathbb R^{n+1}
\]
tal que $\xi(p)\perp T_pM$ para todo $p\in M$.

\end{booktheorem}
\begin{bookproof}{}
\textbf{$(\Leftarrow)$} Supongamos que existe $\xi:M\to\mathbb R^{n+1}$ suave y nunca nula tal que $\xi(p)\perp T_pM$ para todo $p\in M$.

Sea $\omega$ la forma de volumen estandar de $\mathbb R^{n+1}$. Definimos una $n$-forma sobre $M$ por

\[
\eta_p(v_1,\ldots,v_n)=\omega_p(\xi(p),v_1,\ldots,v_n),\qquad v_i\in T_pM.
\]
Esta forma es suave. Ademas no se anula: si $v_1,\ldots,v_n$ es una base de $T_pM$, entonces

\[
\{\xi(p),v_1,\ldots,v_n\}
\]
es base de $\mathbb R^{n+1}$, luego

\[
\omega_p(\xi(p),v_1,\ldots,v_n)\ne0.
\]
Por la proposicion anterior, $M$ es orientable.

\textbf{$(\Rightarrow)$} Supongamos que $M$ es orientable. Queremos construir $\xi$ suave y nunca nulo perpendicular a $T_pM$.

Como $M$ es orientable, existe una $n$-forma suave nunca nula $\eta$ sobre $M$. Localmente, si $(U,\varphi=(x_1,\ldots,x_n))$ es una carta, los vectores

\[
\frac{\partial}{\partial x_1}\bigg|_p,\ldots,\frac{\partial}{\partial x_n}\bigg|_p
\]
forman una base de $T_pM$. En $\mathbb R^{n+1}$ existe exactamente una direccion normal a este subespacio. Elegimos el vector normal unitario $\xi(p)$ de modo que

\[
\omega_p(\xi(p),v_1,\ldots,v_n)
\]
tenga el mismo signo que $\eta_p(v_1,\ldots,v_n)$ para toda base positiva $v_1,\ldots,v_n$ de $T_pM$.

Usando coordenadas locales, esta eleccion se escribe suavemente por cofatores: si

\[
d i_p\left(\frac{\partial}{\partial x_j}\bigg|_p\right)
=\sum_{k=1}^{n+1}a_{kj}(p)\frac{\partial}{\partial y_k},
\]
entonces un vector normal se obtiene como el vector de menores de la matriz $(a_{kj})$. Al normalizar, obtenemos un campo suave porque el vector de menores no se anula.

\bookqed
\end{bookproof}
\begin{bookremark}{}
La orientabilidad de una hipersuperficie se puede pensar como la posibilidad de elegir de manera suave un lado positivo, es decir, un campo normal nunca nulo.

\end{bookremark}
\subsection{Orientacion inducida en el borde}
\begin{bookremark}{}
La orientacion del borde se elige con la convencion "normal saliente primero".

Si $M$ es una variedad orientada con borde y $p\in\partial M$, una base ordenada

\[
(v_1,\ldots,v_{n-1})
\]
de $T_p(\partial M)$ es positiva si

\[
(\nu_p,v_1,\ldots,v_{n-1})
\]
es una base positiva de $T_pM$, donde $\nu_p$ es un vector normal saliente.

\end{bookremark}
\begin{bookexample}{}
En un intervalo $[a,b]$, la orientacion inducida del borde es

\[
\partial[a,b]=\{b\}-\{a\}.
\]
El punto derecho tiene signo positivo y el punto izquierdo tiene signo negativo.

\end{bookexample}
\begin{itemize}
\item [ ]
\end{itemize}
